\documentclass[11pt]{article}

\usepackage[margin=1in]{geometry}
\usepackage{amsmath,amssymb,amsthm}
\usepackage{algorithm}
\usepackage{algpseudocode}
\usepackage{booktabs}
\usepackage{enumitem}
\usepackage{xcolor}
\usepackage{hyperref}

\hypersetup{
  colorlinks=true,
  linkcolor=blue,
  citecolor=blue,
  urlcolor=blue
}

\newtheorem{definition}{Definition}
\newtheorem{theorem}{Theorem}
\newtheorem{lemma}{Lemma}

\newcommand{\Setup}{\mathsf{Setup}}
\newcommand{\Upload}{\mathsf{Upload}}
\newcommand{\Prove}{\mathsf{Prove}}
\newcommand{\Verify}{\mathsf{Verify}}
\newcommand{\KeyGen}{\mathsf{KeyGen}}
\newcommand{\Derive}{\mathsf{Derive}}
\newcommand{\Enc}{\mathsf{Enc}}
\newcommand{\Dec}{\mathsf{Dec}}
\newcommand{\Commit}{\mathsf{Commit}}
\newcommand{\Cover}{\mathsf{Cover}}
\newcommand{\RO}{\mathsf{H}}
\newcommand{\pp}{\mathsf{pp}}
\newcommand{\msk}{\mathsf{msk}}
\newcommand{\sk}{\mathsf{sk}}
\newcommand{\attrs}{\mathsf{attrs}}
\newcommand{\trans}{\mathsf{trans}}

\title{Verifiable Multi-Dimensional Range Queries with M-HIBE Empty-Key Covers and Zero-Knowledge Accumulators}
\author{Anonymous Draft}
\date{\today}

\begin{document}
\maketitle

\begin{abstract}
We study verifiable outsourced multi-dimensional range queries. A server stores a database whose tuples contain $d$ indexed attributes and answers axis-aligned range queries. The client must verify both authenticity, namely that every returned tuple belongs to the outsourced database, and completeness, namely that no matching tuple has been omitted. Existing accumulator-only approaches can prove authenticity succinctly, but their completeness cost grows with the number of excluded points or with a large complement witness. This paper proposes a two-layer construction. Authenticity is proved by a zero-knowledge accumulator over the returned answer set, while completeness is proved by a multi-dimensional HIBE (M-HIBE) empty-key cover. During upload, the data owner issues restricted keys for globally empty prefix-aligned regions. At query time, the server derives query-scoped empty-region keys that cover exactly the empty part of the query box. The verifier checks that the returned answer points and the certified empty regions form a disjoint cover of the query range. We instantiate the method for three-dimensional TPC-H-style queries using WKD-IBE over 12-bit coordinates and a polynomial accumulator. In a 120K-row experiment, the optimized parallel prototype reduces WKD-IBE delegation time from 68.78 seconds to 7.35 seconds and produces 4,543 query-scoped empty keys for a query containing 2,229 occupied spatial points and 24,051 empty spatial points. These preliminary results suggest that M-HIBE empty covers are a practical way to shift completeness proofs from large universal complements to compact, query-local certified empty regions.
\end{abstract}

\section{Introduction}

Outsourced databases are often queried by clients that do not fully trust the storage server. For a multi-dimensional range query
\[
  Q = [\ell_1,u_1] \times \cdots \times [\ell_d,u_d],
\]
the server returns all tuples whose indexed attributes lie in $Q$. A verifiable query protocol should convince the client of two facts. First, every returned tuple is authentic: it was present in the outsourced database. Second, the answer is complete: the server did not omit any tuple satisfying the query predicate.

Authenticity alone can be handled by membership proofs, authenticated data structures, or accumulator-based proofs. Completeness is more subtle. For range queries, the verifier must also be convinced that the region of the query not represented by returned tuples is empty. A direct proof over the full complement of the answer can become expensive in high-dimensional domains, especially when the domain is much larger than the number of matching tuples.

This paper proposes to separate authenticity and completeness. The answer set is handled by a zero-knowledge accumulator proof. The empty part of the query is handled by an M-HIBE empty-key cover. The data owner preprocesses the database and gives the server secret keys for prefix-aligned regions that are known to be empty. A key for a parent empty region can derive keys for subregions, but it cannot be used to derive keys outside that parent. At query time, the server intersects these global empty regions with the query box, derives restricted query-scoped keys, and sends them with the answer. The verifier accepts only if the returned answer points and the empty-region patterns exactly cover the query range.

The resulting proof has a simple geometric interpretation. Let $V_Q$ be the set of grid points inside the query box, let $A_Q$ be the set of unique spatial points represented by the returned answer, and let $E_Q$ be the set of grid points covered by returned empty patterns. Completeness reduces to checking
\[
  A_Q \cup E_Q = V_Q,
  \qquad
  A_Q \cap E_Q = \emptyset,
\]
and checking that each empty pattern is backed by valid M-HIBE key material. If the server omits a real answer point, then that point must either remain uncovered, or be covered by an empty-region key. The first case fails the geometric check; the second would require a valid empty key for a non-empty database region.

\paragraph{Contributions.}
This draft makes the following contributions:
\begin{itemize}[leftmargin=*]
  \item It presents a two-layer verifiable range-query construction combining a ZK accumulator for answer authenticity with M-HIBE empty-region keys for completeness.
  \item It gives a concrete three-dimensional ``XY-parent supplement'' cover that stores empty keys under $X$ prefixes and under $(X,Y)$ prefix parents, avoiding enumeration of the full 3D universe.
  \item It describes a non-interactive query proof in which empty challenges are deterministically derived from the query and answer, and accumulator challenges are derived by Fiat-Shamir transcript hashing.
  \item It reports preliminary implementation results on a TPC-H-style workload, including a serial/parallel ablation for WKD-IBE key delegation.
\end{itemize}

\section{Background}

\subsection{Canonical Range Covers}

Each indexed coordinate is encoded as an $\ell$-bit string. A one-dimensional interval can be decomposed into a disjoint set of canonical binary prefixes. A $d$-dimensional range is then represented as a set of prefix boxes. A prefix box fixes a prefix in each dimension and leaves the remaining suffix bits unrestricted.

For a prefix $p$ of length $|p|$, we write $p \preceq x$ when $p$ is a prefix of the full coordinate string $x$. A wildcard pattern in dimension $i$ has the form
\[
  p_i *^{\ell-|p_i|}.
\]
For $d$ dimensions, a pattern is the concatenation
\[
  P =
  p_1 *^{\ell-|p_1|}
  \parallel \cdots \parallel
  p_d *^{\ell-|p_d|}.
\]
Pattern $P$ contains pattern $P'$ if every fixed bit in $P$ is fixed to the same value in $P'$.

\subsection{WKD-IBE and M-HIBE}

We use a wildcard key-derivation identity-based encryption primitive, denoted WKD-IBE, to implement multi-dimensional HIBE-style keys. The system has public parameters $\pp_{\mathsf{ibe}}$ and a master secret key $\msk_{\mathsf{ibe}}$. A key can be generated for a pattern $P$:
\[
  \sk_P \leftarrow \KeyGen(\msk_{\mathsf{ibe}},P).
\]
If $P$ contains a child pattern $P'$, the holder of $\sk_P$ can derive
\[
  \sk_{P'} \leftarrow \Derive(\sk_P,P').
\]
The security property used by our construction is that a server without a valid ancestor key cannot forge a functional decryption key for an unauthorized pattern, except with negligible probability.

In our setting, M-HIBE keys are not used to decrypt database records. They are used as empty-region certificates. A valid key for pattern $P$ is interpreted as evidence that the data owner certified $P$, or an ancestor of $P$, as empty during preprocessing.

\subsection{Zero-Knowledge Accumulators}

The authenticity layer uses an accumulator or polynomial commitment proof for the returned answer set. Let $D$ be the encoded set of database tuples or spatial points, let $I$ be the encoded returned set, and let $X=D\setminus I$ be the complement. The prover commits to $D$, $X$, and the polynomial representing $I$, then proves that $I$ is consistent with the database digest without revealing additional database elements.

This draft treats the accumulator as an abstract non-interactive proof system
\[
  \pi_{\mathsf{acc}} \leftarrow
  \mathsf{AccProve}(\pp_{\mathsf{acc}},D,I,\trans)
\]
with verifier
\[
  \mathsf{AccVerify}(\pp_{\mathsf{acc}},\mathsf{digest}_D,C_I,\pi_{\mathsf{acc}},\trans).
\]
The concrete implementation uses a bilinear/polynomial accumulator and Fiat-Shamir challenges derived from a transcript.

\section{Problem Definition}

\subsection{System Model}

There are three parties: a data owner, an untrusted server, and a client. The data owner uploads an encoded database $D$ and verification material to the server. Later, a client issues a range query $Q$. The server returns an answer $A$ and a proof object $\Pi$. The client runs $\Verify(Q,A,\Pi)$ and accepts or rejects.

The indexed domain is a grid
\[
  \mathcal{U} = \{0,\ldots,2^\ell-1\}^d.
\]
Multiple database tuples may map to the same spatial point. We therefore distinguish returned rows from unique returned spatial points. The completeness cover is defined over unique spatial points, while tuple-level authenticity is handled by the accumulator layer.

\subsection{Security Goals}

\begin{definition}[Authenticity]
A query proof is authentic if every accepted returned tuple is a member of the outsourced database committed by the data owner.
\end{definition}

\begin{definition}[Completeness]
A query proof is complete if, for every accepted proof for query $Q$, every database tuple whose indexed attributes lie in $Q$ is represented in the returned answer.
\end{definition}

\begin{definition}[Query privacy and leakage]
The proof should not reveal database tuples outside the query answer. The M-HIBE layer reveals query-scoped empty-region patterns and their restricted keys. Thus the intended leakage is the shape of an empty cover inside $Q$, not boundary records or out-of-query tuples.
\end{definition}

\paragraph{Draft scope.}
This draft proves security at the level of a hybrid argument using the security of the accumulator and the unforgeability/functionality of M-HIBE keys. A final version should formalize the exact leakage function and the accumulator relation.

\section{Construction Overview}

The protocol has four phases:
\[
  \Setup,\quad \Upload,\quad \Prove,\quad \Verify.
\]

The data owner first generates M-HIBE and accumulator parameters. During upload, the owner builds prefix occupancy indexes and produces M-HIBE keys for global empty parent regions. These keys are sent to the server together with the outsourced dataset and accumulator material. The server does not receive the M-HIBE master secret key.

When a query arrives, the server computes the answer set and the query-scoped empty cover. For every selected empty pattern $E$, the server finds a global empty parent $R$ such that $R \supseteq E$ and derives $\sk_E$ from $\sk_R$. It then produces a ZK accumulator proof for the returned answer. The proof object is
\[
  \Pi =
  \left(
    \mathcal{E}_Q,
    \{\sk_E\}_{E\in\mathcal{E}_Q},
    \pi_{\mathsf{acc}}
  \right).
\]

The client verifies the accumulator proof, verifies the geometry of the returned empty patterns, and checks the validity of the M-HIBE keys for their claimed patterns.

\section{M-HIBE Empty-Key Cover}

\subsection{Multi-Dimensional Pattern Encoding}

For $d$ dimensions and $\ell$ bits per coordinate, every pattern has length $d\ell$. The first $\ell$ slots encode dimension 1, the next $\ell$ slots encode dimension 2, and so on. In the three-dimensional implementation,
\[
  P = p_x *^{\ell-|p_x|}
  \parallel
  p_y *^{\ell-|p_y|}
  \parallel
  p_z *^{\ell-|p_z|}.
\]
The attribute list for WKD-IBE contains one attribute for every fixed bit. Wildcard slots are omitted.

\subsection{3D XY-Parent Supplement Cover}

The implemented construction specializes the general prefix idea to three dimensions:
\[
  X=\textsf{shipdate},\qquad
  Y=\textsf{encoded discount},\qquad
  Z=\textsf{quantity}.
\]
The prototype uses $\ell=12$, hence 36 WKD-IBE slots.

During upload, the owner builds two occupancy maps:
\[
  \mathsf{XToY}[p_x]
  =
  \{y : \exists (x,y,z)\in D,\ p_x\preceq x\},
\]
\[
  \mathsf{XYToZ}[p_x,p_y]
  =
  \{z : \exists (x,y,z)\in D,\ p_x\preceq x,\ p_y\preceq y\}.
\]
For each $X$ prefix, the owner finds empty $Y$ prefixes and creates patterns
\[
  p_x \parallel e_y \parallel *^\ell.
\]
For each occupied $(p_x,p_y)$ parent, the owner finds empty $Z$ prefixes and creates patterns
\[
  p_x \parallel p_y \parallel e_z.
\]
Redundant subregions are removed by keeping maximal patterns.

This construction avoids materializing every empty 3D box. It stores parent empty regions aligned with $X$ and $(X,Y)$ prefix nodes, and the server derives query-scoped children only when a query arrives.

\begin{algorithm}[t]
\caption{Upload-time empty parent key generation}
\label{alg:upload-empty}
\begin{algorithmic}[1]
\Require Dataset $D$, bit length $\ell$, M-HIBE master secret key $\msk_{\mathsf{ibe}}$
\Ensure Global empty parent key set $\mathcal{K}$
\State Build $\mathsf{XToY}$ and $\mathsf{XYToZ}$ from all points in $D$
\State $\mathcal{R} \gets \emptyset$
\For{each indexed $X$ prefix $p_x$}
  \State $\mathcal{Y}_{\mathsf{empty}} \gets \mathsf{CanonicalEmptyCover}([0,2^\ell-1]\setminus \mathsf{XToY}[p_x])$
  \For{each $e_y\in \mathcal{Y}_{\mathsf{empty}}$}
    \State $\mathcal{R}\gets \mathcal{R}\cup\{p_x\parallel e_y\parallel *^\ell\}$
  \EndFor
\EndFor
\For{each indexed pair $(p_x,p_y)$}
  \State $\mathcal{Z}_{\mathsf{empty}} \gets \mathsf{CanonicalEmptyCover}([0,2^\ell-1]\setminus \mathsf{XYToZ}[p_x,p_y])$
  \For{each $e_z\in \mathcal{Z}_{\mathsf{empty}}$}
    \State $\mathcal{R}\gets \mathcal{R}\cup\{p_x\parallel p_y\parallel e_z\}$
  \EndFor
\EndFor
\State Remove non-maximal patterns from $\mathcal{R}$
\State $\mathcal{K}\gets \{(R,\KeyGen(\msk_{\mathsf{ibe}},R)) : R\in \mathcal{R}\}$
\State \Return $\mathcal{K}$
\end{algorithmic}
\end{algorithm}

\subsection{Query-Scoped Empty Patterns}

The global parent regions may extend outside the query. Therefore the server intersects each touched global empty parent with $Q$, decomposes every non-empty intersection back into canonical prefix boxes, and then minimizes overlaps over the empty points in $Q$.

Let $A_Q$ be the unique returned spatial points. The target empty space is
\[
  V_Q \setminus A_Q.
\]
The selected empty patterns $\mathcal{E}_Q$ should satisfy
\[
  \bigcup_{E\in\mathcal{E}_Q} E = V_Q\setminus A_Q,
  \qquad
  E\cap A_Q=\emptyset \text{ for all } E\in\mathcal{E}_Q.
\]

\begin{algorithm}[t]
\caption{Server proof generation}
\label{alg:prove}
\begin{algorithmic}[1]
\Require Query $Q$, database $D$, global empty parent keys $\mathcal{K}$, accumulator material
\Ensure Answer $A$ and proof $\Pi$
\State $A\gets \{r\in D: r \text{ satisfies } Q\}$
\State $A_Q\gets$ unique spatial points represented by $A$
\State $\mathcal{C}\gets\emptyset$
\For{each $(R,\sk_R)\in\mathcal{K}$ touched by $Q$}
  \State $I\gets R\cap Q$
  \If{$I\neq \emptyset$}
    \State $\mathcal{C}\gets \mathcal{C}\cup \mathsf{CanonicalCover}(I)$
  \EndIf
\EndFor
\State $\mathcal{E}_Q\gets \mathsf{GreedySetCover}(\mathcal{C},V_Q\setminus A_Q)$
\For{each $E\in\mathcal{E}_Q$}
  \State Find $(R,\sk_R)\in\mathcal{K}$ such that $R\supseteq E$
  \State $\sk_E\gets \Derive(\sk_R,E)$
\EndFor
\State $\trans\gets \RO(\pp,Q,\mathsf{Hash}(A),\mathsf{Hash}(\mathcal{E}_Q))$
\State $\pi_{\mathsf{acc}}\gets \mathsf{AccProve}(\pp_{\mathsf{acc}},D,A,\trans)$
\State $\Pi\gets(\mathcal{E}_Q,\{\sk_E\}_{E\in\mathcal{E}_Q},\pi_{\mathsf{acc}})$
\State \Return $(A,\Pi)$
\end{algorithmic}
\end{algorithm}

\section{Verification}

Verification has three parts. First, the verifier checks answer authenticity with the accumulator proof. Second, it checks the geometry of the empty cover. Third, it validates the M-HIBE key material.

\begin{algorithm}[t]
\caption{Client verification}
\label{alg:verify}
\begin{algorithmic}[1]
\Require Query $Q$, answer $A$, proof $\Pi=(\mathcal{E}_Q,\{\sk_E\},\pi_{\mathsf{acc}})$
\Ensure Accept or reject
\State $A_Q\gets$ unique spatial points represented by $A$
\State $\trans\gets \RO(\pp,Q,\mathsf{Hash}(A),\mathsf{Hash}(\mathcal{E}_Q))$
\If{$\mathsf{AccVerify}(\pp_{\mathsf{acc}},\mathsf{digest}_D,A,\pi_{\mathsf{acc}},\trans)=0$}
  \State \Return reject
\EndIf
\For{each $E\in\mathcal{E}_Q$}
  \If{$E\not\subseteq Q$ or $E\cap A_Q\neq \emptyset$}
    \State \Return reject
  \EndIf
\EndFor
\If{$A_Q\cup \bigcup_{E\in\mathcal{E}_Q}E \neq V_Q$}
  \State \Return reject
\EndIf
\For{each checked $E\in\mathcal{E}_Q$}
  \If{$\mathsf{CheckKey}(\pp_{\mathsf{ibe}},E,\sk_E)=0$}
    \State \Return reject
  \EndIf
\EndFor
\State \Return accept
\end{algorithmic}
\end{algorithm}

\paragraph{Functional key check.}
The prototype checks a claimed key by encrypting a fresh random message under the claimed pattern and decrypting it with the supplied key:
\[
\begin{aligned}
\mathsf{CheckKey}(\pp_{\mathsf{ibe}},P,\sk_P):\quad
&\attrs \leftarrow \mathsf{PatternToAttributes}(P),\\
&m \leftarrow \mathsf{RandomMessage}(),\\
&ct \leftarrow \Enc(\pp_{\mathsf{ibe}},\attrs,m),\\
&m' \leftarrow \Dec(ct,\sk_P),\\
&\mathsf{return}\; [m'=m].
\end{aligned}
\]
A malformed or unrelated key fails except with negligible probability. The benchmark samples at most 64 keys for this check. A full verification mode can check every returned empty key with cost linear in $|\mathcal{E}_Q|$.

\section{Non-Interactive Proof}

The original challenge view is interactive: after seeing the answer, the client computes empty challenge regions and asks the server to prove them. We make this non-interactive in two steps.

First, the empty challenge set is deterministic. It is derived from the query and answer geometry:
\[
  \mathcal{E}_Q = \Cover(Q,A_Q).
\]
In the optimized implementation, the server proposes the cover by intersecting global empty parents with the query and minimizing the resulting query-scoped patterns. The verifier checks the resulting geometry directly.

Second, accumulator randomness is derived using Fiat-Shamir. The transcript should bind the public parameters, dataset digest, query bounds, answer commitment or answer hash, empty-pattern list hash, and previous proof messages. A representative challenge has the form
\[
  \rho =
  \RO(\pp,Q,C_A,\mathsf{digest}_D,\mathsf{Hash}(\mathcal{E}_Q),\mathsf{proof\_prefix}).
\]
The server sends one proof message after receiving $Q$:
\[
  \mathsf{Server}\rightarrow\mathsf{Client}:
  \left(A,\mathcal{E}_Q,\{\sk_E\}_{E\in\mathcal{E}_Q},\pi_{\mathsf{acc}}\right).
\]

\section{Security Argument}

\begin{theorem}[Authenticity, informal]
If the accumulator proof system is sound, then a polynomial-time server cannot make the verifier accept a returned tuple that is not in the outsourced database, except with negligible probability.
\end{theorem}

\begin{proof}[Proof sketch]
The verifier accepts only if $\mathsf{AccVerify}$ accepts for the returned answer and the transcript binding the query and proof. A forged or substituted tuple would imply an accepting accumulator proof for a false statement. This contradicts accumulator soundness.
\end{proof}

\begin{theorem}[Completeness, informal]
Assume the data owner issues M-HIBE keys only for empty regions and that the M-HIBE key-derivation scheme is unforgeable for unauthorized patterns. Then a server cannot omit a matching spatial point and still pass full verification, except with negligible probability.
\end{theorem}

\begin{proof}[Proof sketch]
Suppose a matching point $p\in V_Q$ is omitted from the returned answer. Since the geometric check requires the answer points and empty regions to cover $V_Q$, point $p$ must be covered by some empty pattern $E\in\mathcal{E}_Q$. Since $p$ is a real database point, $E$ is not an empty region. The server can pass the key check for $E$ only if it has a valid M-HIBE key for $E$ or for an ancestor certified by the owner. But the owner issued keys only for empty regions, so such an ancestor cannot contain $p$. Therefore the server must forge unauthorized M-HIBE key material, contradicting M-HIBE security.
\end{proof}

\begin{theorem}[Zero-knowledge leakage, informal]
Assuming the accumulator proof is zero-knowledge and the only revealed M-HIBE material is query-scoped empty keys for patterns inside $Q$, the proof reveals no database tuple outside the returned answer beyond the empty-cover leakage inside the query.
\end{theorem}

\begin{proof}[Proof sketch]
The accumulator layer is zero-knowledge by assumption. The M-HIBE layer reveals the list of empty patterns and their restricted keys. These patterns are contained in $Q$ and are checked to avoid returned answer points. Therefore the explicit leakage is the shape of the empty cover within $Q$. No out-of-query tuple or boundary record is returned by the protocol. A final proof should formalize a simulator that receives this leakage and simulates the proof object.
\end{proof}

\section{Implementation}

The prototype is implemented in Go on top of a WKD-IBE library using BLS12-381. The three query dimensions are shipdate, encoded discount, and quantity. Each coordinate is encoded using $\ell=12$ bits, so a 3D pattern has 36 WKD-IBE slots.

The implementation has two logical engines. Engine A computes the M-HIBE completeness proof: canonical query keys, global parent intersections, query-scoped empty patterns, parent-key lookup, and child-key derivation. Engine B computes the accumulator proof for the returned answer.

\paragraph{Parallelization.}
The parallel implementation does not change the selected cover or the proof format. It parallelizes independent operations: finding a parent for each scoped pattern, generating unique parent keys in the benchmark simulation, deriving child keys grouped by parent, and sampling key-decryption checks during verification. With one worker, it behaves like the serial version.

\paragraph{Prototype caveat.}
In the intended deployment, the server stores owner-generated global empty parent keys:
\[
  \mathsf{Owner}: \msk_{\mathsf{ibe}}
  \rightarrow
  \{\sk_R\}_{R\in\mathcal{R}}
  \rightarrow
  \mathsf{Server}.
\]
The benchmark derives touched parent keys from the master key inside a combined program to simulate the offline materialized key material and measure how much parent key material would be touched. The paper protocol should describe the real deployment, not the benchmark shortcut.

\section{Evaluation}

This section reports preliminary measurements. The main 3D query is a TPC-H Q6-like range query over 120K rows. The implementation uses 20 worker goroutines for the optimized parallel M-HIBE prover. Reported times are wall-clock times.

\begin{table}[t]
\centering
\caption{Logical proof shape for the 120K-row 3D query.}
\label{tab:proof-shape}
\begin{tabular}{lr}
\toprule
Metric & Value \\
\midrule
Query canonical range keys & 24 \\
Global empty parent regions touched & 20,985 \\
Query-scoped empty regions selected & 4,543 \\
Matching rows & 2,327 \\
Unique matching spatial points & 2,229 \\
Empty spatial points covered & 24,051 \\
Total query spatial volume & 26,280 \\
\bottomrule
\end{tabular}
\end{table}

\begin{table}[t]
\centering
\caption{Serial and parallel 3D performance on the 120K-row query.}
\label{tab:parallel}
\begin{tabular}{lrr}
\toprule
Metric & Serial & Parallel, 20 workers \\
\midrule
Engine A total (s) & 77.62 & 16.32 \\
WKD-IBE delegation (s) & 68.78 & 7.35 \\
Client completeness check (s) & 1.65 & 0.16 \\
Total server proving (s) & 100.96 & 39.96 \\
\bottomrule
\end{tabular}
\end{table}

\begin{table}[t]
\centering
\caption{Proof and key material sizes for the 120K-row 3D query.}
\label{tab:sizes}
\begin{tabular}{lr}
\toprule
Material & Size \\
\midrule
Query range key material & 14.27 KB \\
Global parent key material touched & 1,737.30 KB \\
Query-scoped empty key material & 1,553.73 KB \\
ZK proof size & 1.20 KB \\
\bottomrule
\end{tabular}
\end{table}

Parallelization affects only independent M-HIBE key derivations and verifier sampling checks. It does not change the selected empty-region cover, the number of returned empty keys, or the accumulator proof. Thus we report the parallel wall-clock time as the optimized result and the serial time as an ablation isolating the benefit of multi-core key derivation.

\begin{table}[t]
\centering
\caption{Preliminary baseline comparison. PoneglyphDB entries are not yet measured under the same 3D setting.}
\label{tab:baseline}
\begin{tabular}{lrrr}
\toprule
Scheme & Server time (s) & Client verify (ms) & VO size (KB) \\
\midrule
Pure ZK accumulator & 7.91 & 5.38 & 2.52 \\
Ours, M-HIBE+ZK & 39.96 & 160.16 & 1,569.20 \\
PoneglyphDB & N/A & N/A & N/A \\
\bottomrule
\end{tabular}
\end{table}

The pure accumulator baseline is compact and fast in this preliminary table, but it is measured in a trapdoor or full-completeness baseline mode and should be described carefully. The current M-HIBE proof has larger verification-object size because it transmits explicit query-scoped empty keys. Its advantage is that completeness is certified by query-local empty regions rather than by a large universal complement proof. A final comparison requires all baselines to use the same hardware, query, security level, and leakage assumptions.

\section{Limitations and Open Issues}

\paragraph{Owner-generated parent keys.}
The final paper must state clearly that the server stores owner-generated global empty parent keys. The benchmark currently derives touched parent keys from the master key inside the experiment driver. This is a simulation of offline materialized keys, not the real trust model.

\paragraph{Sampled key checks.}
The reported benchmark samples at most 64 derived empty keys for same-pattern encryption/decryption checks. The geometric coverage check is exhaustive over the query grid, but the key-functional check is sampled. A full-key audit experiment should be added or the sampled check should be labeled as a performance-oriented sanity check.

\paragraph{Accumulator transcript.}
The non-interactive ZK section should define the Fiat-Shamir transcript exactly. It should bind public parameters, dataset digest, query bounds, answer commitment, empty-pattern list hash, accumulator messages, and any implementation-specific domain separators.

\paragraph{Verifier knowledge in the benchmark.}
The benchmark runs in a combined program where the true dataset is available, so it can compute the real matching spatial volume directly. The protocol description must instead define $A_Q$ from the returned answer and rely on the accumulator proof for authenticity.

\paragraph{Revealing keys versus proving knowledge.}
The current implementation reveals restricted query-scoped empty keys and verifies them by functional decryption tests. An alternative is to prove knowledge of those keys. Revealing keys is simpler and matches the prototype, but the leakage statement should be written around explicit empty-key disclosure.

\paragraph{Evaluation completeness.}
Dataset-size scalability, query-selectivity sensitivity, PoneglyphDB comparison under the same setting, and full-key audit overhead remain to be measured before a final empirical claim.

\section{Related Work}

This draft should be completed with precise citations for authenticated range queries, accumulators, zero-knowledge set proofs, HIBE/WKD-IBE, and verifiable databases such as PoneglyphDB. For now, the text uses these primitives abstractly and focuses on the new combination of accumulator authenticity with M-HIBE empty-region completeness.

\section{Conclusion}

We presented a draft construction for verifiable multi-dimensional range queries that separates authenticity from completeness. A ZK accumulator proves that returned tuples belong to the outsourced database, while M-HIBE empty-region keys prove that the rest of the query box is empty. The three-dimensional XY-parent supplement cover avoids enumerating the full grid and supports query-time derivation of restricted empty keys. Preliminary results show that parallel WKD-IBE delegation substantially improves proving time, while the main remaining costs are explicit empty-key material and accumulator computation. The next steps are to formalize leakage and transcripts, add full-key audit measurements, and complete baseline comparisons under a unified experimental setup.

\bibliographystyle{plain}
\begin{thebibliography}{9}

\bibitem{wkdibe}
TODO: Add the WKD-IBE reference used by the implementation.

\bibitem{accumulators}
TODO: Add the polynomial or bilinear accumulator reference used by the ZK layer.

\bibitem{poneglyphdb}
TODO: Add the PoneglyphDB reference and clarify the exact baseline settings.

\bibitem{tpch}
TODO: Add the TPC-H benchmark reference.

\end{thebibliography}

\end{document}
